\documentclass[10pt]{article}
\usepackage[margin=0.75in]{geometry}
\usepackage{booktabs}
\usepackage{microtype}
\usepackage{enumitem}
\usepackage{xcolor}
\usepackage{hyperref}
\hypersetup{colorlinks=true, urlcolor=blue, linkcolor=black}
\usepackage{titlesec}
\titlespacing*{\section}{0pt}{6pt}{2pt}
\titlespacing*{\subsection}{0pt}{4pt}{1pt}
\setlist{nosep, topsep=2pt, partopsep=0pt, leftmargin=*}
\setlength{\parindent}{0pt}
\setlength{\parskip}{4pt}

\title{\vspace{-0.6in}\Large CSE/DSC 234 Project 2 -- Schema Linking with SFT\\
\large Final Report}
\author{}
\date{}

\begin{document}
\maketitle\vspace{-0.4in}

\textbf{Final validation leaderboard: 0.7270} (Table 0.7911 / Column 0.6630).
Baseline (smoke adapter): 0.5605. Best single adapter: 0.6648 (sweep 13).
Submission: \texttt{python3 main.py --input X --output Y} runs a \textbf{4-way
LoRA ensemble with majority-2 aggregation} \emph{across two base models AND
mixed retrievers per adapter}: sweep 22 on Qwen3-1.7B + embed, sweep 13 on
Qwen-Coder-1.5B + embed, sweep 15 on Qwen-Coder-1.5B + BM25, and sweep 13
\emph{reused} on Qwen-Coder-1.5B with hybrid retrieval. An identifier
survives iff it is emitted by $\geq$2 of the 4 passes. Completes 101
questions in \textbf{13m9s} on a 24~GB MIG slice ($<$15~min budget).
\texttt{main.py} reads each adapter's \texttt{adapter\_config.json} to group
adapters by base model (so each base is loaded only once) and accepts a
\texttt{path[:retriever]} spec per adapter, so the same adapter under
different retrievers can be ensembled too. The 4th pass is exactly that --
pure post-hoc diversity, no extra training.

\section{Training Data Methodology}

The released 301 train + 101 validation split was the starting point. We treated
schema linking as next-token prediction over a chat prompt: a fixed system message
explains the JSON contract and wildcard convention; the user turn carries the
serialized schema and NL question; the assistant turn carries the canonical JSON
target (sorted keys + sorted columns) emitted by \texttt{prompt.target\_string}.

\textbf{Schema serialization.} The baseline serializer is one line per table,
\texttt{TableName: c1, c2, \dots}. We tested two richer variants:
(i)~\textbf{types} -- appending column type in parentheses; (ii)~\textbf{keys} --
\texttt{*} for primary keys and \texttt{->TargetTable.col} for foreign keys.
A combined \texttt{types\_keys} serialization roughly doubled token counts on the
dense SBO schemas; on the largest (SBO-Banking) the prompt exceeded 6k tokens.
\texttt{keys} alone added no tokens on SBO because the SBO Spider schemas carry
no PK/FK metadata -- only 5 of the 17 schemas do (ATBI, NTSB,
NorthernPlainsFireManagement, KlamathInvasiveSpecies, PacificIslandLandbirds).

\textbf{Output format.} Raw JSON with stable casing, alphabetical keys and column
lists. A table with no specific columns (wildcard queries like \texttt{count(*) from
t}) is emitted as \texttt{\{"t": []\}}, not omitted. We verified this end-to-end
with \texttt{tests/test\_checklist.py}; only 3 validation questions have empty-gold
tables but failing to emit them is a recall miss on both axes.

\textbf{Table-level retrieval at training time.} Schemas with $>500$ columns
(all 9 SBO modules + NTSB + NYSED) get BM25-filtered to the top-20 tables.
\textbf{Oracle augmentation}: at training time we force-include every
gold-referenced table in the kept set so the model is never asked to predict
labels absent from its prompt. Inference does not have this oracle.

\textbf{Augmentation.} The 9 SBO modules ship with only 6--12 train examples each
($\sim$22\% of train) versus 24--60 for the parks/NTSB domains; SBO is also where
$\sim$55\% of our validation table-misses concentrate. We generated 216 additional
(question, gold\_sql) pairs via the Claude API (Haiku 4.5; 9 batches of 24)
focused entirely on SBO, labeled each with the project's
\texttt{data\_prep/sql\_to\_schema\_links.py} (sqlglot-based) extractor, and dropped any
that failed to parse or referenced 0 schema tables. Only 1 of 217 candidates was
dropped (a duplicate question; 0 hallucinated tables, 0 parse errors -- evidence
the prompt and validation pipeline are tight). The merged
\texttt{train\_augmented.json} has 517 examples. The audit log lives at
\texttt{data/train\_augmented\_audit.json}.

\section{Final Pipeline Architecture}

\textbf{Base model.} \texttt{Qwen/Qwen2.5-Coder-1.5B-Instruct}. The code-pretrained
variant of Qwen2.5-1.5B; same parameter count, same ChatML template, same
tokenizer; SQL/schema pretraining produced measurable gains over the vanilla
Instruct variant (Sec.~4).

\textbf{Adapter.} LoRA r=32, $\alpha=64$, dropout 0.05, target modules
\{\texttt{q,k,v,o,gate,up,down}\}\_\texttt{proj} (the full attention block plus
the MLP block). bf16 with gradient checkpointing during training. Adapter
safetensors are shipped as \textbf{fp16} (PEFT's default fp32 save was 148~MB
per adapter, above GitHub's 100~MB per-file limit; fp16 halves that to 74~MB
and we verified the score is preserved within $\pm 0.001$ of fp32). The 3
adapters together fit in $\sim$220 MB of repo size.

\textbf{Optimizer.} AdamW (TRL default), \texttt{lr\_scheduler\_type=linear},
\texttt{warmup\_ratio=0.05}, \texttt{per\_device\_train\_batch\_size=1},
\texttt{gradient\_accumulation\_steps=4}, \texttt{max\_steps=200},
\texttt{max\_length=4096}. Three adapters are shipped: sweep 13 (Qwen-Coder-1.5B,
lr=3e-4, 301 train), sweep 15 (Qwen-Coder-1.5B, lr=3e-4 + augmented data,
400 steps), and sweep 22 (Qwen3-1.7B, lr=3e-4, 301 train). The 4th inference
pass reuses sweep 13's weights with a different retriever -- so the 4-way
ensemble runs over 3 trained adapters, not 4.

\textbf{Inference pipeline (\texttt{main.py}).} Adapters are grouped by their
declared base model (read from \texttt{adapter\_config.json}); each base loads
once. Within a group, adapters are loaded as named PEFT adapters and swapped
via \texttt{model.set\_adapter}. For each (adapter, retriever) pass we
accumulate per-question $\{$table$\rightarrow$cols$\}$ predictions and per-pass
vote counts; after the last pass we apply \texttt{--aggregation maj2}
(default; \texttt{union} also supported). \textbf{Partial-save safety
contract}: passes 1--3 write the running union, pass 4 writes \texttt{maj2}.
Combined with the deliberate adapter order (Qwen3 first), this means: after
3 passes the on-disk file matches the previous locked champion (0.7107); a
mid-pass-4 kill ships that baseline; clean completion ships 0.7270.
\textbf{Retrieval}: \texttt{BAAI/bge-small-en-v1.5} cosine similarity (top-20
tables) for passes 1--2, BM25 for pass 3 (sweep 15), RRF hybrid for pass 4
(sweep 13). The cross-retriever diversity is what \texttt{maj2} converts
into score.

\textbf{Post-processing.} \texttt{canonicalize\_prediction} drops any
non-schema identifier (precision-preserving hallucination filter) and re-cases
identifiers to schema casing. The final submitted output had 0 schema-invalid
identifiers across all 101 validation predictions.

\section{Experimentation Methodology}

We ran 20 numbered sweeps via \textbf{RapidFire AI}
(\texttt{training/train\_rapidfire.py}, \texttt{RFGridSearch}/\texttt{RFSFTConfig}). Sweeps
1--9 were spent surviving infrastructure problems (described below); sweeps
10--20 produced the substantive results in Table~\ref{tab:configs}. Per-sweep
\texttt{rapidfire.log}, \texttt{training.log}, and extracted MLflow metrics
(\texttt{metrics.json}) live under \texttt{logs/sweepN/}; full per-run configs
are pulled from the rapidfireai MLflow store on demand.

\subsection{Sweep dimensions}
Across the three \emph{training-time} axes -- base model, adapter capacity,
training data -- and four \emph{inference-time} axes -- prompt style, retriever,
ensemble strategy, post-processing -- we ran 11 distinct training configs and
$\sim$10 distinct inference-time variants. Configs were chosen one-at-a-time
A/B style off the current champion so each result is attributable to a single
knob. Multi-config sweeps (8--16 in \texttt{RFList}) were used in our first
attempts (sweeps 5--8) but, after infrastructure issues forced us to
\texttt{num\_chunks=1} (Sec.~3.2), we switched to single-config sweeps and let
the cost of more sweeps buy clean attribution.

\subsection{Infrastructure issues and how we worked around them}

\textbf{Issue \#1: DSMLP cleanup SIGKILLs ($\sim$24h of lost GPU time).}
Sweeps 5--8 ran multi-config \texttt{RFGridSearch} with 8--16 configs and
\texttt{max\_steps}=300. Each tried to train through $\sim$5--6h of wall time.
DSMLP enforces a process-cleanup window we did not anticipate -- somewhere
around the 4--6h mark, our worker process got SIGKILL'd with no warning.
The mlflow metrics survived (they are written to a SQLite DB on
\texttt{\$HOME}) so we could see that sweep 8 had reached step 264/300
across all 8 configs before dying. But the trained LoRA weights themselves had
\emph{nowhere on disk}. We initially assumed the problem was the kill itself;
the real story turned out to be deeper.

\textbf{Issue \#2: RapidFire's SHM-only checkpoint store.} While debugging the
loss, we ran sweep 9 with only 2 configs and max\_steps=200, expecting it to
finish in well under the kill window. It did -- ``Experiment ended'' fired
cleanly, no SIGKILL. \emph{We still had no
\texttt{final\_checkpoint/*.safetensors} on disk.} Reading
\texttt{rapidfireai/fit/backend/worker.py} around L440 revealed why. By default
RapidFire writes intermediate checkpoints to \texttt{/dev/shm} via its
\texttt{shm\_manager}; the disk write only fires inside a single branch gated on
\texttt{is\_run\_finished == (chunk\_id == num\_chunks-1) \&\&
(steps $\geq$ total\_steps)}. With the controller's round-robin chunk schedule,
a run typically hits its step budget on chunk $N$ where $N <$ num\_chunks$-1$,
the disk-save branch is skipped, and the in-SHM adapter dies with the worker
process at \texttt{experiment.end()}. We considered patching the constant
\texttt{USE\_SHARED\_MEMORY = True} at module level but rejected that as
fragile (it affects all rapidfireai usage on the system and would not survive
a \texttt{pip install --upgrade}). The clean workaround: pass
\texttt{--num\_chunks 1} so every run necessarily ends on ``the last chunk'' by
construction and the disk-save branch fires reliably. All sweeps from 10 onward
use this workaround; sweep 10 was the first successful end-to-end persisted
RapidFire training in our series.

\textbf{Issue \#3: HuggingFace Trainer fallback (\texttt{training/train\_single.py}).}
While debugging Issue \#2, we were not sure whether the problem was orchestration
(rapidfireai) or training (TRL/PEFT setup). We built
\texttt{training/train\_single.py} -- a $\sim$120-line plain
\texttt{trl.SFTTrainer} script -- as an A/B against the same config. It
trained fine in 14 minutes and produced a 0.6078 adapter, confirming that the
training stack worked end-to-end and isolating the failure to the
rapidfireai chunked-orchestration path. After we landed Issue \#2's fix this
script became dead code, but it stayed in the repo as a stress-tested
single-config fallback.

\textbf{Issue \#4: CUDA SIGFPE in \texttt{main.py} at \texttt{batch\_size>1}.}
On the first $\sim$0.61 inference run we hit a hard ``Floating point exception
(core dumped)'' right at the first \texttt{model.generate} call. Process
crashed before producing any predictions. We narrowed the trigger by bisecting
the CLI flags and found that batch size 1 worked, $\geq$2 SIGFPE'd. Stack
combination is \texttt{transformers 4.57 + torch 2.10 + bf16 + left-padded
prompts $\sim$4k tokens + H100}. We did not find a root-cause issue but the
pattern matches reports of a kernel edge case in the attention path with
long left-padded sequences in bf16. We defaulted
\texttt{batch\_size=1} in \texttt{main.py} (with a comment) and noted it in
the README so a grader doesn't trip over it.

\textbf{Issue \#5: Mismatched val-data layout.} Halfway through development we
reorganized the repo to push raw data under \texttt{data/}; the
\texttt{tests/test\_checklist.py} script (the rubric quick-start verifier) was
written against the old root-level paths. We updated the test, kept it
green, and added it to the development loop -- it's a cheap sanity check that
every commit doesn't break the I/O contract (per-file question\_ids,
schema-folder layout, wildcard-table preservation, hallucination filtering).

\subsection{The supporting tools we built}
\texttt{data\_prep/} holds data-side scripts: \texttt{augment\_data.py}
(SBO augmentation via Haiku 4.5, SQL validated through
\texttt{sql\_to\_schema\_links.py}, audit log preserved),
\texttt{format\_training\_data.py} / \texttt{format\_two\_stage\_data.py}
(JSON $\rightarrow$ JSONL), \texttt{generate\_cot\_traces.py} (sweep 21
rationales), and \texttt{expand\_columns.py} (Sec.~4 post-processor).
\texttt{scripts/probe\_ensembles.py} is the combinatorial probe that
discovered the 4-way + \texttt{maj2} (138 cells, calls \texttt{eval.py}
on each, ranks by score). \texttt{logs/README.md} maps each sweep to
its config and metrics from the rapidfireai MLflow store.

\begin{table}[t]
\centering\small
\setlength\tabcolsep{4pt}
\begin{tabular}{rllrrrr}
\toprule
\# & base model & key knobs & lr & steps & val LB & $\Delta$ \\
\midrule
10 & Qwen2.5-1.5B          & r=16 q/k/v/o, BM25                                & 3e-4 & 200 & 0.6229 & --- \\
12 & Qwen2.5-\textbf{Coder}-1.5B & r=16 q/k/v/o, BM25                          & 3e-4 & 200 & 0.6258 & +0.003 \\
\textbf{13} & Qwen2.5-Coder-1.5B & \textbf{r=32 + MLP}, BM25                   & 3e-4 & 200 & \textbf{0.6648} & \textbf{+0.039} \\
18 & Qwen2.5-Coder-1.5B    & r=32 + MLP, BM25                                  & 2e-4 & 200 & 0.6609 & $-$0.004 \\
19 & Qwen2.5-Coder-1.5B    & r=32 + MLP, BM25, \textbf{cosine LR}              & 3e-4 & 200 & 0.6479 & $-$0.017 \\
14 & Qwen2.5-Coder-1.5B    & r=32 + MLP, BM25, \textbf{517 train (aug)}        & 3e-4 & 200 & 0.6379 & $-$0.027 \\
15 & Qwen2.5-Coder-1.5B    & r=32 + MLP, BM25, 517 train, \textbf{400 steps}   & 3e-4 & 400 & 0.6600 & $-$0.005 \\
16 & Qwen2.5-Coder-1.5B    & r=32 + MLP, BM25, 301 train, \textbf{400 steps}   & 3e-4 & 400 & 0.6268 & $-$0.038 \\
11 & Qwen2.5-1.5B          & r=16 q/k/v/o, \textbf{types\_keys prompt}         & 3e-4 & 200 & 0.5946 & $-$0.028 \\
20 & \textbf{SmolLM2-1.7B} & r=32 + MLP, BM25                                  & 3e-4 & 200 & 0.5301 & $-$0.135 \\
17a/b & Qwen2.5-Coder-1.5B & \textbf{Two-stage}: table model + column model    & 3e-4 & 200 & 0.6127 & $-$0.052 \\
21 & Qwen2.5-Coder-1.5B    & r=32 + MLP, \textbf{CoT-distilled} training data  & 3e-4 & 200 & 0.6103 & $-$0.054 \\
22 & \textbf{Qwen3-1.7B}   & r=32 + MLP, BM25, original data                   & 3e-4 & 200 & 0.6599 & $-$0.049 \\
27 & Qwen2.5-Coder-1.5B    & r=32 + MLP, \textbf{embed-prompted train}         & 3e-4 & 200 & 0.6772 & $-$0.001 vs s13 embed \\
28 & Qwen2.5-Coder-1.5B    & r=32 + MLP, \textbf{keys prompt (PK/FK markers)}  & 3e-4 & 200 & 0.6464 & $-$0.032 \\
\midrule
\multicolumn{6}{r}{\textbf{Inference-time -- best single adapter (sweep 13):}} & \\
\multicolumn{6}{r}{BM25 retrieval (baseline)} & 0.6648 \\
\multicolumn{6}{r}{max\_tables: 20 $\rightarrow$ 40 (more distractors)} & 0.6032 \\
\multicolumn{6}{r}{\textbf{Embedding retrieval (BAAI/bge-small-en-v1.5)}} & \textbf{0.6786} \\
\multicolumn{6}{r}{Hybrid retrieval (RRF of BM25 + embed)} & 0.6724 \\
\multicolumn{6}{r}{Self-consistency, K=3, T=0.5, union-of-samples} & 0.6737 \\
\multicolumn{6}{r}{Column expansion (question-keyword $\leftrightarrow$ col-name match)} & $\sim$same \\
\midrule
\multicolumn{6}{r}{\textbf{Ensembles -- 3-way Coder-only (sweeps 13+15+18):}} & \\
\multicolumn{6}{r}{BM25 retrieval} & 0.6986 \\
\multicolumn{6}{r}{Embedding retrieval (single-base champion before Qwen3)}    & 0.7002 \\
\multicolumn{6}{r}{Hybrid retrieval (RRF)}                                    & 0.6977 \\
\multicolumn{6}{r}{2-way 13+18, 4-way +14 / +20 / +11}                        & 0.66--0.69 \\
\multicolumn{6}{r}{\textbf{Ensembles -- cross-family (Coder + Qwen3):}} & \\
\multicolumn{6}{r}{2-way 13+22 (both embed)} & 0.7063 \\
\multicolumn{6}{r}{3-way 13+15+22 all-embed (single-retriever)} & 0.7019 \\
\multicolumn{6}{r}{3-way mixed-retriever: 13\_embed + 15\_BM25 + 22\_embed (previous champion)}
                                                              & 0.7107 \\
\multicolumn{6}{r}{4-way + \texttt{maj2}: above three + 13\_hybrid, $\geq$2-vote aggregation} & \\
\multicolumn{6}{r}{\textbf{(submitted champion: same adapters, no extra training; 13m9s wall time)}}
                                                              & \textbf{0.7270} \\
\multicolumn{6}{r}{$~~~$same 4-way under partial-union aggregation (regression -- shows \texttt{maj2} is doing the work)} & 0.7034 \\
\bottomrule
\end{tabular}
\caption{All 13 distinct training configs (top) and the inference-time ablations
on the best single adapter and the 3-way ensemble (bottom). Bold marks the
single best config in each section. $\Delta$ is vs sweep~10 (the
single-config baseline). \texttt{r=32 + MLP} means LoRA r=32 with target
modules \{q,k,v,o,gate,up,down\}\_proj. \texttt{BM25} retrieval scores tables
by query-vs-(table+cols) BM25; \texttt{embed} replaces BM25 with sentence-
transformer cosine similarity (\texttt{BAAI/bge-small-en-v1.5}); \texttt{hybrid}
is Reciprocal Rank Fusion of the two.}
\label{tab:configs}
\end{table}

\subsection{Inference-time experimentation}
Once we had the single-best adapter (sweep~13), we tested eight inference-time
variants: max\_tables 20$\rightarrow$40, three retrievers (BM25 / embed /
hybrid via RRF), self-consistency at K=3 (union-of-samples and
majority-vote), column-expansion post-processing, union/intersection/
majority-k ensembling over progressively more adapters (2--6 way), and
\emph{mixed-retriever passes over the same adapter} (sweep 13 once with
embed and once with hybrid as separate ensemble members).
Variants were A/B'd directly against the previous champion file in
\texttt{predictions/} -- no retraining, $\sim$5--15 min per A/B; the
4-way + maj2 search was a combinatorial offline probe over 138 (combo,
aggregation) cells against the per-(adapter, retriever) prediction files
we had accumulated. Many variants were negative (Sec.~4) but two stood out:
embedding retrieval (+0.014 single, +0.016 in the ensemble), and the
4-way + \texttt{maj2} combination (+0.016 over the 3-way mixed-retriever).

\section{Results and Analysis}

\textbf{What worked.}
\begin{itemize}
\item \textbf{Switching to Coder-1.5B (+0.003)}: small but in the right direction;
sets up the bigger gains.
\item \textbf{Bigger LoRA: r=32 + MLP modules (+0.039)}: the single biggest
training-time lever. The full attention+MLP adapter has $\sim$3.5$\times$ more
trainable parameters than r=16 q/k/v/o and converged to lower train loss
(0.075 vs 0.108) without overfitting at 200 steps. Notably this is the
configuration our sweeps 5--7 attempted but lost to the SHM bug -- the
\texttt{num\_chunks=1} workaround was load-bearing for actually \emph{getting}
this number.
\item \textbf{Embedding-based table retrieval (+0.014 single, +0.016 ensemble)}:
\texttt{BAAI/bge-small-en-v1.5} surfaces semantically related tables that BM25
misses (e.g.\ qid 95 needs \texttt{ORCT}, ranked outside BM25 top-10 because
all top hits are \texttt{PWZ}* payment-wizard children; the embedder puts it
at rank 1). The biggest \emph{inference-time} lever.
\item \textbf{Union ensemble of 3 diverse adapters (+0.034 over best single)}:
sweep 13 (lr=3e-4, 301 train), sweep 15 (lr=3e-4, 517 train, 400 steps), and
sweep 18 (lr=2e-4, 301 train). Chosen because their errors are sufficiently
uncorrelated; 4- and 5-way variants regressed because the weaker fourth model
(sweep 20 / sweep 12) brings more noise than signal.
\item \textbf{Adding Qwen3-1.7B (sweep 22) for cross-family ensemble (+0.017
over single-base champion)}: sweep 22 alone scores 0.66 -- weaker than
sweep 13 (0.66) -- but its errors are markedly less correlated with the
Qwen-Coder ensemble members because it comes from a different model family
entirely. Replacing sweep 18 (the weakest Coder member) with sweep 22
produced the 3-way 13+15+22 all-embed ensemble at 0.7019. The submission's
\texttt{main.py} groups adapters by base model and loads each base only
once, so this multi-base ensemble still fits in the 15-minute inference
budget.
\item \textbf{Mixed-retriever ensemble (+0.009 over the all-embed champion)}:
the same adapter produces different predictions under different retrievers
because the candidate-table set in the prompt differs. Putting sweep 15
under BM25 (while 13 and 22 stay on embed) gave us 0.7107 -- the right
retriever for each ensemble member is whichever decorrelates its errors
from the others, not the one used at training time. Exposed in
\texttt{main.py}'s \texttt{--ensemble\_dirs} via a \texttt{path[:retriever]}
spec per adapter.
\item \textbf{4-way + majority-2 aggregation (+0.016)}: a combinatorial
offline probe over our accumulated per-(adapter, retriever) prediction
files found that adding a 4th pass -- sweep 13 reused under hybrid
retrieval, no extra training -- and switching aggregation from union to
\textbf{majority-2} ($\geq$2 of 4 votes) lifts the leaderboard to
\textbf{0.7270}. The same 4-way under union actually regresses to 0.7034:
\texttt{maj2} is doing the work, not the extra pass. Mechanically this is
the precision/recall trade we needed -- union maximizes recall (table F1
was already 0.79); \texttt{maj2} prunes column-level false positives. Both
axes improved (column 0.6423$\rightarrow$0.6630, table 0.7791$\rightarrow$0.7911).
Wall time 13m9s, inside the 15-min budget.
\end{itemize}

\textbf{What didn't work, and why.}
\begin{itemize}
\item \textbf{\texttt{types\_keys} prompt (-0.028)}: roughly doubled token count
and \emph{regressed} both table and column scores. We attribute this to attention
dilution -- with $\sim$2$\times$ more schema text per prompt, the 1.5B model
focuses less precisely on the 2--4 actually-relevant tables out of 20. Cheap
signal at this scale beats expensive signal.
\item \textbf{\texttt{max\_tables}=40 (-0.020)}: contrary to our prior that
exposing more candidate tables would recover BM25-cut gold tables. Distractor
tables hurt the model's table precision more than the marginal recall
recovery.
\item \textbf{Two-stage (table model then column model) (-0.052)}: a clean
architectural negative result. Stage-A loses the column signal that
disambiguates similarly-named SBO siblings (\texttt{OPMG} vs \texttt{OPHA});
stage-B cannot reason cross-table about JOIN columns; errors compound (a wrong
table in stage-A can't be corrected by stage-B). Joint reasoning is load-bearing
for this task.
\item \textbf{SmolLM2-1.7B base swap (sweep 20, -0.135 alone)}: the model is
genuinely weaker on this distribution (token accuracy 0.80 vs Coder's 0.90);
adding it to the ensemble \emph{hurt} (-0.037), so diversity does not always
help when the noise floor of the new member is too high.
\item \textbf{Data augmentation alone (sweep 14, -0.027)}: train\_loss
$\sim$2$\times$ higher than the 301-example baseline at 200 steps
(underfitting); doubling steps (sweep 15) recovered most of the regression
but did not exceed the baseline. The 2$\times$2 table-size/step-count
grid implies the augmented data is roughly neutral here -- the limiting
resource is per-example exposure, not raw data quantity.
\item \textbf{Self-consistency union-of-samples (-0.005)} and
\textbf{column-expansion post-processing (neutral)}: T=0.5 sampling
adds noise; question-keyword $\leftrightarrow$ col-name match either
over-fires (\textgreater 1 token: $-$0.009) or under-fires
(\textgreater 2 meaningful tokens: only $\sim$9 cols, wash). Greedy
decode + the precision-preserving canonicaliser do the work already.
\item \textbf{Chain-of-thought distillation (sweep 21, -0.054)}: brief
($\sim$50 token) Haiku-generated rationales were schema-grounded
(300/301 referenced at least one gold table) but the CoT-trained adapter
landed at $\sim$0.61 and regressed the ensemble by 0.007. Likely a mix
of: gradient budget spent on reasoning tokens, descriptive rather than
disambiguating rationales, and early commitment that propagates errors
into the JSON.
\item \textbf{Embed-prompted training (sweep 27, neutral)}: sweep 13's
recipe with the training-time table filter switched from BM25 to embed.
The prior was that closing the small train/test retrieval gap would
yield +0.005--0.02; standalone the adapter scores 0.6772 (vs sweep 13's
0.6786 under embed inference -- identical), and swapping it into the
4-way lifts LB +0.002 entirely on the table side (column F1
0.6630$\rightarrow$0.6631). Retrieval mismatch is not the column-F1
bottleneck.
\item \textbf{Keys prompt style (sweep 28, -0.032 alone, regresses every
ensemble swap by 0.008--0.028)}: trained with PK \texttt{*} markers and
\texttt{->TargetTable.col} FK arrows, expecting it to disambiguate
FK-shared columns (\texttt{CASEID} across NTSB tables -- the top
column-F1 FP class in our error analysis). The FK metadata annotates
\emph{boring} columns (\texttt{VERSION->CHILDSEAT.VERSION}); the
actually-confusing ones (\texttt{CASEID}, \texttt{EVENT\_ID},
\texttt{VEHNO}) appear without markers because the schema doesn't list
them as FKs. The prompt added structure for non-problems.
\end{itemize}

\textbf{Which knobs mattered for table vs column F1.}
The submitted 4-way's table score (0.7911) is 13 points above its column
score (0.6630), and this asymmetry is consistent across all our adapters
and across the iterative champions (3-way locked: 0.7648 / 0.6356; 4-way
maj2: 0.7911 / 0.6630). Sweep~13's r=32+MLP upgrade lifted column F1 the
most among training changes (+0.045 vs sweep~12); embedding retrieval
lifted column recall (+0.149 over BM25 in the 3-way ensemble). The
\texttt{maj2} switch was the first inference-time lever that improved
\emph{both} axes (table +0.012, column +0.021): it shaves the union's
column-level false positives without losing the recall that drove the
3-way's table side. Going the other way, the SmolLM2 swap collapsed both
axes evenly -- suggesting column performance is bottlenecked by base-model
capacity and retrieval quality, not by prompt or post-processing.

\textbf{What we would try next.}
(1)~A column-specialist second stage conditioned on the 4-way's table
predictions (which are already 0.79); sweep 17's two-stage failed
because stage-A had nontrivial table noise -- conditioning on
ensemble-agreed tables changes that failure mode. (2)~A bigger base for
cross-family ensembling (Coder-3B); sweeps 27 and 28 suggest the
column-F1 ceiling around 0.66 is set by base capacity, not prompt or
data. (3)~Hybrid retrieval with a learned cross-encoder reranker over
BM25 $\cup$ embed top-30 candidates.

\section{AI Tools Statement}

Code was developed iteratively with Anthropic's \textbf{Claude Code}
(\textbf{Claude Opus 4.7}). Specific contributions: drafting sweep
configurations, investigating the rapidfireai SHM-checkpoint disk-flush bug
(reading \texttt{rapidfireai/fit/backend/worker.py} to identify the
\texttt{num\_chunks=1} workaround), implementing the embedding retriever,
the LoRA-swap ensemble path in \texttt{main.py} (\texttt{--aggregation maj2},
the partial-write safety contract, the ordered ensemble spec), the
combinatorial probe that discovered the 4-way + \texttt{maj2} lift to
0.7270, the post-processing and reproduction utilities, the SBO augmentation
pipeline, and the CoT-trace generator -- both of which call the Anthropic
API using \textbf{Claude Haiku 4.5} -- and drafting this report. All
design decisions, sweep selection, eval interpretation, and final
configuration choices were ours; the AI tool served as an implementation
and analysis accelerant. The audit log
\texttt{data/train\_augmented\_audit.json} records every generated example
and its validation outcome for transparency.

\end{document}
